\begin{cnabstract}
摘要的撰写应符合《文献摘要编写规则》（GB/T 6447）的规定，中文摘要字数一般为300字至1000字。摘要应概括研究目的、方法、结果和结论，突出研究成果。摘要中尽量避免采用图、表、化学结构式、非公知公用的符号和术语。

第一段主要阐述选题的意义，说明研究问题的背景和价值。

第二段简单介绍研究动机、方法和目的。从存在的问题出发，阐述采用的方法、依据的理论或进行的实验，说明解决的问题及取得的成果。本文的主要工作和成果如下：

1. 针对具体研究问题，说明采用的方法、提出的方案及验证结果。所有结论应有真实研究数据支持。

2. 摘要应表达学位论文的中心内容，简短明了，摘取论文中的基本信息，体现科研工作的核心思想，根据内容自行分段描述。

最后一段对全文进行总结，并对进一步的研究提出展望。中英文摘要内容应基本对应，关键词应选择公知公用的词或术语，不宜采用自造词及过长的短语。
\cnkeywords{材质检索；材质优化；三维重建；深度学习}
\end{cnabstract}
